\documentclass[protokoll]{dvd}

\KOMAoptions{
    headwidth = 18cm,
    footwidth = 18cm,
}

% Styrelsen
\def\ordf{Daniell Cole}
\def\vice{\vak}
\def\samo{Angelika Hagberg}
\def\kassa{Owais Tabbaa}
\def\sekr{\vak}
\def\vak{\emph{vakant}}


% mötesinfo

\def\nr{6}
\def\yr{2026}
\def\ymd{\yr-11-17}
\def\starttime{12:00}


\begin{document}

\title{Styrelsemöte \nr}
\subtitle{\yr}
\author{Styrelsen}
\date{\ymd}


\textbf{Datum:} \csname @date\endcsname\\
\textbf{Tid:} \starttime\\
\textbf{Plats:} Styrelserummet\\
\textbf{Styrelsemedlemmar:}
\begin{närvarande_förtroendevalda}
    \förtroendevald{Ordförande}{\ordf}{Ja}
    \förtroendevald{Vice Ordförande}{\vice}{---}
    \förtroendevald{Kassör}{\kassa}{Ja}
    \förtroendevald{SAMO}{\samo}{Ja}
    \förtroendevald{Sekreterare}{\sekr}{---}
\end{närvarande_förtroendevalda}

%\textbf{Övriga medlemmar:}
%\begin{närvarande_medlemmar}
%\medlem{Ingen}[]
%\end{närvarande_medlemmar}
\section{Öppnande av möte}

Mötet beräknas öppnas av \ordf\ kl \starttime

Mötet öppnades av \ordf\ kl 12.08


\section{Runda bordet}

Runda bordet innebär att varje person berättar hur de känner sig.
Man kan till exempel berätta att man är stressad på grund av en inlämning, irriterad på sin granne, eller bara väldigt glad därför att man ligger i fas med plugget.

\section{Formalia}

\subsection{Styrelsens beslutbarhet}

\blockquote[7 kap. 5 \S~första stycket i stadgan][]{%
    Styrelsen är endast beslutsmässig då samtliga styrelsemedlemmar har fått kallelsen till styrelsemötet och minst hälften av styrelsemedlemmarna är närvarande. Ordförande eller vice ordförande måste vara närvarande när beslut tas.
}

\blockquote[7 kap. 11 \S~första stycket i stadgan]{
Styrelsens sammanträden måste kallas minst två läsdagar innan sammanträdet. }

\blockquote[7 kap. 11 \S~sista stycket i stadgan]{
Kallelse till sammanträde får endast ske när protokollet från tidigare sammanträde är justerat och tillgängligt för medlemmarna.}

\subsubsection*{Förslag}

\begin{attsatser}
    \item Styrelsen har uppnått kraven i 7 kap. 5 § första stycket i stadgan och är därmed beslutbar.
\end{attsatser}
\subsubsection*{Beslut}
\begin{attsatser}
    \item %förslaget till beslut bifalles
\end{attsatser}


\subsection{Fastställande av mötesschema}

För att styrelsen ska kunna fatta ett styrelsebeslut eller protokollföra en diskussion behöver punkten i mötesschemat där styrelsen ska fatta beslut vara inlagd eller föras in i mötesschemat senast vid den här punkten.

\subsubsection*{Förslag}

\begin{attsatser}
    \item 
    mötesschemat fastställs utan några förändringar.
    %punkt läggs till på mötesschema
\end{attsatser}
\subsubsection*{Beslut}
\begin{attsatser}
    \item förslaget till beslut bifalles
\end{attsatser}

\subsection{Val av mötessekreterare}
Owais volunterar.

\subsubsection*{Förslag}
\begin{attsatser}
    \item Owais väljs till mötessekreterare. 
\end{attsatser}
\subsubsection*{Beslut}
\begin{attsatser}
    \item förslaget till beslut bifalles
\end{attsatser}

\subsection{Val av protokolljusterare}

Protokolljusterare har till uppgift att kontrollera att protokollet i slutändan reflekterar de faktiska besluten och diskussionerna som fördes under mötet.
Utöver protokolljusteraren så ska mötesordförande och mötessekreteraren signera protokollet.
Vid styrelsemöten ska det endast vara en justerare.
Mötesordförande och mötessekreteraren kan inte vara justerare.

\subsubsection*{Förslag}
\begin{attsatser}
    \item Angelika väjls till protokolljusterare.
\end{attsatser}
\subsubsection*{Beslut}
\begin{attsatser}
    \item förslaget till beslut bifalles
\end{attsatser}

\section{Rapport} 
\subsection{Ordförande}
Snackat med Albin om stadgaändringar, hade bra synpunkter, tar det under disskussionspunkt. 
Jag glömde boka sal till stämma, ska fixa det idag.

Skickade ut notis om att strep sökes, bara Owais svarat.

\subsection{Kassör}
Ska bokföra alla äskningar för detta året + Albins vattenkokare, inte frågat
Hannes hur man gör årsrapport (frågar snart) sen ska kolla på förra årets äskningar också. 

\subsection{SAMO}
Inget att rapportera

\newpage


\section{Diskussionspunkter}
\subsection*{Feedback från Albin}
\subsubsection*{Ändringsförslag sammanfattning}
Bättre om det står i ett ''guide''-dokument alt på hemsidan som sammanfattar medlemskapsrättigheter och refererar till stadgan och dess relevanta avsnitt. 

Stadgan var skriven med tanken att varje kapitel för sig definerade sin sak utan något överlap. Denna ändring skulle riskera att läsaren skulle behöva hoppa mellan kapitel. 

Det finns också en risk som denna ändring skulle kunna medföra: Om en framida ändring av stadgan skulle ske, finns det en risk att den ändringen \textit{bara} sker i sitt relevanta kapitel och detta inte också uppadteras i paragrafen om medlemars rättigheter, då kan det ske att stadgan står emot sig själv. 


\subsubsection*{Ändringsförslag styrelseledamot}
Riskerar att urhålka styrelsen och i framtida år riskerar att alltid ha minimal styrelse men med flertal uppdragslösa medlemmar. 

Bättre att försöka främja en kutlur av att kommitté-representanter deltar på styrelsemöten, kanske inte nödvändigtvist hela möten utan en del i början, t.ex. tom rapporter?

Alternativt definiera en mer specifik roll, t.ex. kommittérepresentant via vissa kommittéer har positioner i styrelsen. 

Ett annat alternativ är att vi försöker hålla mer frekventa kommittémöten. 

\section{Besultspunkter}

\subsection{Stadgaändringar}
\subsubsection*{Ändringsförslag 1}
Minimum 3 medlemmar i styrelsen: 
\begin{attsatser}
    \item i början av 7 kap. 3 § i stadgan införa 
    \begin{quote}
        Styrelsen kan som minst bestå av tre medlemmar varav en skall vara divisionsordförande och en divisionskassör. 
    \end{quote}
\end{attsatser}

\subsubsection*{Ändringsförslag 2}
Krav från Göta, styrelsemedlem behöver vara ordinarie medlem i Göta studentkår:
\begin{attsatser}
    \item ändra första stycket i 7 kap. 2 § från
    \begin{quote}
        Styrelsens medlemmar väljs av divisionsstämman. Mandatperioden är densamma som räkenskapsåret. Valbar till styrelsen är medlem i divisionen.
    \end{quote}
    till
    \begin{quote}
        Styrelsens medlemmar väljs av divisionsstämman. Mandatperioden är densamma som räkenskapsåret. Valbar till styrelsen är medlem i divisionen med ordinarie medlemskap i Göta studentkår.
    \end{quote}
\end{attsatser}

\subsubsection{Förslag till beslut}
\begin{attsatser}
    \item lägga till stadgaändringsförslag 1 \& 2
\end{attsatser}
\subsubsection{Beslut}
\begin{attsatser}
    \item Attsats 1 bifalles.
\end{attsatser}

\subsection{Kommittémöte}

\begin{attsatser}
    \item Owais skickar ut om möte
\end{attsatser}
\subsubsection{Beslut}
\begin{attsatser}
    \item %Attsats 1 bifalles.
\end{attsatser}

\subsection{Äskan-utdlykt}
\begin{quote}
Hej!

Vi I DVRK26 vill arrangera en utflykt tillsammans med de kommande DVRK27. Vi vill göra detta för att öka gemenskapen inom DVRK, och ha en tydlig övergång till nästa generation. Syftet med utflykten är att vi ska kunna ha lära-känna aktiviteter tillsammans, där alla de olika rollerna kan få prata mer ingående kring sin roll, och vad den nästa generationen kan förvänta sig. 

Vi vill hyra en stuga under en natt, där vi kan samlas, umgås och leka roliga lekar, äta god mat, och så småningom spendera natten där. Tanken är att pengarna ska gå till stughyrning, mat, snacks och eventuella förnödenheter. 

Vi vill äska för 4 000 kr. 

Vänligen, 
Ola "A'LOCO" Forss
\end{quote}
\begin{attsatser}
    \item Fint initiativ men... Nej. Vi kan inte godkänna en äskan för ett slutet event. 
\end{attsatser}
\subsubsection{Beslut} 
\begin{attsatser}     
\item Äskan avslås. 
\end{attsatser}


\section{Per Capsulam-beslut}
\blockquote[]{
När per capsulam-beslut har fattats måste beslutet dokumenteras,\\
\ \cdot\ \ på divisionens Discordserver som är tillgänglig alla studenter;\\
\ \cdot\ \ i protokollet för det första styrelsesammanträdet efter beslut har fattats.\\

Dokumentationen ska innehålla\\
\ \cdot\ \ vilka att-satser som har bifallits;\\
\ \cdot\ \ omröstningsperioden.\\

Om dessa krav inte uppfylls upphävs beslutet.
}
\textit{Inga Per Capsulam-beslut har fattats.}


\newpage
\section{Avslutande av möte}

\subsection{Nästa möte}
2026-12-01 - 12:00

\subsection{Mötets avslutande}
Mötet avslutades kl 13:08 

\styrelsesignaturer

\end{document}